\documentclass[12pt,letterpaper]{article}

\usepackage[utf8]{inputenc}
\usepackage[T1]{fontenc}
\usepackage{amsmath}
\usepackage{amsfonts}
\usepackage{amssymb}
\usepackage{bm}
\usepackage{graphicx}
\usepackage{natbib}
\usepackage{url}
\usepackage{hyperref}
\usepackage{geometry}
\usepackage{booktabs}
\usepackage{array}
\usepackage{float}

\geometry{margin=1in}
\setlength{\parskip}{6pt}

\renewcommand{\thesection}{S\arabic{section}}
\renewcommand{\thefigure}{S\arabic{figure}}
\renewcommand{\thetable}{S\arabic{table}}
\renewcommand{\theequation}{S\arabic{equation}}

\title{Supplementary Information:\\
  Synchronization confounds inference of coupling in dynamical systems}
\date{}

\begin{document}
\maketitle

\section{Sensitivity Equations}\label{sec:si_sens}

Computing the Cram\'er-Rao lower bound and optimizing the likelihood
both require derivatives of model outputs with respect to parameters.
Let $\phi \in \{\theta, S_1(0), S_2(0), I_1(0), I_2(0)\}$ denote a
parameter ($\phi = \phi_k$ for some $k = 1, \dots, 5$). Denote
$\mathbf{S}(t) = (S_1(t), S_2(t))$, $\mathbf{I}(t) = (I_1(t),
I_2(t))$, and $\mathbf{N}^{-1} = \text{diag}(1/N_1, 1/N_2)$.

The sensitivity of the force of infection is:
\begin{equation}\label{eq:si_sens_lambda}
\frac{\partial \lambda_i(t)}{\partial \phi} = \beta_i(t) \left(
\frac{\partial \mathbf{C}}{\partial \phi} \mathbf{N}^{-1}
\mathbf{I}(t) + \mathbf{C} \mathbf{N}^{-1} \frac{\partial
  \mathbf{I}(t)}{\partial \phi} \right)_i
\end{equation}
where
$\frac{\partial \mathbf{C}}{\partial \theta} =
\begin{bmatrix} -1 & 1 \\ 1 & -1 \end{bmatrix}$
and $\partial \mathbf{C}/\partial \phi = \mathbf{0}$ for $\phi \neq
\theta$.

The sensitivity of the expected incidence follows by the product rule:
\begin{equation}\label{eq:si_sens_mu}
\frac{\partial \mu_i(t)}{\partial \phi} =
\frac{\partial S_i(t)}{\partial \phi} [1 - \exp(-\lambda_i(t))]
+ S_i(t) \exp(-\lambda_i(t)) \frac{\partial \lambda_i(t)}{\partial \phi}
\end{equation}

The state sensitivities propagate according to:
\begin{align}\begin{split}\label{eq:si_sens}
\frac{\partial S_i(t+1)}{\partial \phi} &= \exp(-\lambda_i(t))
\frac{\partial S_i(t)}{\partial \phi} - S_i(t) \exp(-\lambda_i(t))
\frac{\partial \lambda_i(t)}{\partial \phi} \\
\frac{\partial I_i(t+1)}{\partial \phi} &= \exp(-\gamma)
\frac{\partial I_i(t)}{\partial \phi} +
\frac{\partial \mu_i(t)}{\partial \phi}
\end{split}\end{align}
with initial conditions $\partial S_i(0)/\partial S_j(0) =
\delta_{ij}$, $\partial I_i(0)/\partial I_j(0) = \delta_{ij}$, and
zero otherwise. These are solved jointly with the state equations in a
single forward pass.


\section{Derivation of the Fisher Information Matrix}\label{sec:si_fim}

Differentiating the log-likelihood with respect to $\phi$ and
defining $A_i(t) := -\frac{1}{2\mu_i(t)} + \frac{Y_i(t) -
  \rho\mu_i(t)}{(1-\rho)\mu_i(t)} + \frac{(Y_i(t) -
  \rho\mu_i(t))^2}{2\rho(1-\rho)\mu_i(t)^2}$, we obtain $\partial
\ell / \partial \phi = \sum_{t,i} A_i(t)\, \partial
\mu_i(t)/\partial \phi$. Since $\mathbb{E}[A_i(t)] = 0$ and
observations are independent across $(i,t)$:
\begin{equation*}
J_{kl} = \sum_{t,i} \mathbb{E}[A_i(t)^2]\,
\frac{\partial \mu_i(t)}{\partial \phi_k}
\frac{\partial \mu_i(t)}{\partial \phi_l}
\end{equation*}

Writing $\epsilon_i(t) = Y_i(t) - \rho\mu_i(t) \sim
\mathcal{N}(0,\rho(1-\rho)\mu_i(t))$ and expanding $A_i(t)^2$:
\begin{align*}
A_i(t)^2 &= \frac{1}{4\mu_i(t)^2} -
\frac{\epsilon_i(t)}{(1-\rho)\mu_i(t)^2} -
\frac{\epsilon_i(t)^2}{2\rho(1-\rho)\mu_i(t)^3} \\
&\quad + \frac{\epsilon_i(t)^2}{(1-\rho)^2\mu_i(t)^2} +
\frac{\epsilon_i(t)^3}{\rho(1-\rho)^2\mu_i(t)^3} +
\frac{\epsilon_i(t)^4}{4\rho^2(1-\rho)^2\mu_i(t)^4}
\end{align*}

Using Gaussian moments ($\mathbb{E}[\epsilon^2] =
\rho(1-\rho)\mu_i(t)$, $\mathbb{E}[\epsilon^3] = 0$,
$\mathbb{E}[\epsilon^4] = 3[\rho(1-\rho)\mu_i(t)]^2$):
\begin{align*}
\mathbb{E}[A_i(t)^2] &= \frac{1}{4\mu_i(t)^2} -
\frac{1}{2\mu_i(t)^2} + \frac{\rho}{(1-\rho)\mu_i(t)^2} +
\frac{3}{4\mu_i(t)^2} = \frac{1+\rho}{2(1-\rho)\mu_i(t)^2}
\end{align*}

Therefore:
\begin{equation}\label{eq:si_fim_scalar}
J_{kl} = \frac{1+\rho}{2(1-\rho)}\sum_{t=0}^{T-1}\sum_{i=1}^{2}
\frac{1}{\mu_i(t)^2} \frac{\partial \mu_i(t)}{\partial \phi_k}
\frac{\partial \mu_i(t)}{\partial \phi_l}
\end{equation}
In matrix form, $\mathbf{J} = \mathbf{G}^T \mathbf{W} \mathbf{G}$
where $\mathbf{G} \in \mathbb{R}^{2T \times 5}$ has entries
$G_{(i,t),k} = \partial \mu_i(t)/\partial \phi_k$ and $\mathbf{W} =
\frac{1+\rho}{2(1-\rho)}\,\text{diag}(\mu_1(0)^{-2}, \ldots,
\mu_2(T{-}1)^{-2})$.


\section{Model Parameters}\label{sec:si_params}

\begin{table}[h]
\centering
\begin{tabular}{llll}
\toprule
Parameter & Value & Description & Source or calculation\\
\midrule
$R^{\text{max}}_0$ & 3.52 & $R_0$ at peak season & \cite{shaman2010} \\
$R^{\text{min}}_0$ & 1.12 & $R_0$ during the summer & \cite{shaman2010} \\
$D$ & 3.24 & Mean infection period (days) & \cite{shaman2010} \\
$\gamma$ & $2.16$ & Weekly recovery rate & $7 / D$ \\
$\beta_{\text{max}}$ & $3.11$ & Peak-season transmission & $R_0^{\text{max}}(1 - e^{-\gamma})$\\
$\beta_{\text{min}}$ & $0.99$ & Low-season transmission & $R_0^{\text{min}}(1 - e^{-\gamma})$ \\
$\beta_0$ & $2.05$ & Mean transmission rate & $(\beta_{\text{min}} + \beta_{\text{max}})/2$ \\
$\delta$ & $0.517$ & Seasonal amplitude & $\beta_{\text{max}} / \beta_0 - 1$\\
$\rho$ & $5.36 \times 10^{-4}$ & Infection Fatality Rate & \citet{rolfes2018}\\
$\varphi$ & & Phase of seasonal forcing & per state \\
\bottomrule
\end{tabular}
\caption{Model parameters calibrated to seasonal influenza with
  $\beta(t) = \beta_0(1+\delta \sin(2\pi t + \varphi))$.}
\label{tab:si_params}
\end{table}


\section{Statistical Power Heatmap}\label{sec:si_power}

Figure~\ref{fig:si_power} shows the statistical power for detecting
non-zero connectivity as a function of the true connectivity and
seasonal amplitude.

\begin{figure}[h]
\centering
\includegraphics[width=\textwidth]{../pix/power.pdf}
\caption{Statistical power for detecting $\theta > 0$ at the 5\%
  significance level. Left: synchronized (in-phase). Right:
  asynchronous (out-of-phase by $\pi$). When epidemics are
  synchronized, power remains low even for large $\theta$.}
\label{fig:si_power}
\end{figure}


\section{Data and Reconstruction}\label{sec:si_data}

\begin{figure}[h]
\centering
\includegraphics[width=\textwidth]{../pix/raw_deaths.pdf}
\caption{Weekly P\&I death counts for California, New York, Texas, and
  Florida, showing seasonal peaks superimposed on a year-round
  non-influenza baseline.}
\label{fig:si_raw}
\end{figure}

\begin{figure}[h]
\centering
\includegraphics[width=\textwidth]{../pix/seasons_highlighted.pdf}
\caption{P\&I deaths with flu seasons used for inference highlighted
  (darker segments). COVID-19 pandemic years (2020--2022) are
  excluded.}
\label{fig:si_seasons}
\end{figure}

\begin{figure}[h]
\centering
\includegraphics[width=\textwidth]{../pix/rec_CAxNY.pdf}
\caption{Reconstruction of excess P\&I deaths for California (top) and
  New York (bottom) for three flu seasons. Black dots: observed weekly
  excess deaths. Red line: fitted model $\rho\mu_i(t)$.}
\label{fig:si_reconstruction}
\end{figure}


\bibliographystyle{plainnat}
\bibliography{refs.bib, ~/lib.bib}

\end{document}
